\section{The Frameworks across the Moments: An Analytic Matrix}
\label{sec:matrix}

\noindent
The frameworks of \xref{sec:bg-five} are drawn upon at the moments of the cycle to which each is suited, and the paper's treatment of the moments (\xref{sec:cognition} through \xref{sec:recognition}) develops this drawing-upon in prose. Because the relation of five frameworks to five moments is not easily held in view across the sections that develop it, this section sets it out in a matrix, which serves as a map of the treatment to follow and as a statement, in compact form, of the foundational claim that each framework contributes a distinct epistemology and a distinct practice at each moment. The matrix distinguishes two layers, since each framework, at each moment, both cognises the moment's object in a characteristic way, its epistemological layer, and guides the moment's action in a characteristic way, its practical layer. The two layers are presented in two tables, the first for the epistemology and the second for the practice.

The matrix is to be read as an analytic device and not as a partition. The moments interpenetrate (\xref{sec:cyc-moments}), and the frameworks are not assigned exclusively to moments but are each available at each moment, more central at some and more peripheral at others; the entries record the characteristic contribution of each framework at each moment, the contribution the prose develops where it is most consequential. A cell records how a given framework, at a given moment, takes up that moment's task, and the reader tracing a single framework down a column will find its distinctive standpoint maintained across the moments, while the reader tracing a single moment across a row will find the polyphony of frameworks that the moment draws upon.

\begin{table}[t]
\centering
\scriptsize
\renewcommand{\arraystretch}{1.3}
\caption{The epistemological layer: how each framework cognises the object of each moment of the cycle.}
\label{tab:matrix-epistemic}
\begin{tabularx}{\linewidth}{@{}>{\raggedright\arraybackslash}p{1.7cm} >{\raggedright\arraybackslash}X >{\raggedright\arraybackslash}X >{\raggedright\arraybackslash}X >{\raggedright\arraybackslash}X >{\raggedright\arraybackslash}X@{}}
\toprule
\rowcolor{pageblush}
\textbf{\textcolor{rosedeep}{Framework}} & \textbf{\textcolor{rosedeep}{Cognition}} & \textbf{\textcolor{rosedeep}{Practice}} & \textbf{\textcolor{rosedeep}{Evaluation}} & \textbf{\textcolor{rosedeep}{Reflection}} & \textbf{\textcolor{rosedeep}{Re-cognition}} \\
\midrule
Positivist & Flourishing as a measurable matter of fact & Practices as interventions with measurable effects & Measurement by validated instruments & Measures read against their limits & Revision of measured baselines \\
Evidentialist & Belief about the good proportioned to evidence & Supported practices, warranted by controlled evidence & Belief updating on the evidence of the lived fabric & Re-proportioning of belief as evidence revises & Carrying forward the updated belief \\
Normative & Flourishing as the substantive freedom of each to unfold & Practice as the realisation of capability & Assessment against the specification of just flourishing & Whether the specification was met, and at whose cost & Re-specification of the good in light of the turn \\
Virtue-theoretic & Flourishing as activity in accordance with cultivated excellence & Habituation; formation of \emph{hexis} and family manner & The state of the dispositions formed & Whether the dispositions are excellences or their contraries & Disposition carried forward, deepened or narrowed \\
Phenomenological & Flourishing as lived and interpreted before measured & Cultivation of attention to the shared life & Thick description; the felt and meant quality & The lived meaning the measures discarded & The reopened, unsymbolised apprehension \\
\bottomrule
\end{tabularx}
\end{table}

\begin{table}[t]
\centering
\scriptsize
\renewcommand{\arraystretch}{1.3}
\caption{The practical layer: how each framework guides the action of each moment of the cycle.}
\label{tab:matrix-practical}
\begin{tabularx}{\linewidth}{@{}>{\raggedright\arraybackslash}p{1.7cm} >{\raggedright\arraybackslash}X >{\raggedright\arraybackslash}X >{\raggedright\arraybackslash}X >{\raggedright\arraybackslash}X >{\raggedright\arraybackslash}X@{}}
\toprule
\rowcolor{pageblush}
\textbf{\textcolor{rosedeep}{Framework}} & \textbf{\textcolor{rosedeep}{Cognition}} & \textbf{\textcolor{rosedeep}{Practice}} & \textbf{\textcolor{rosedeep}{Evaluation}} & \textbf{\textcolor{rosedeep}{Reflection}} & \textbf{\textcolor{rosedeep}{Re-cognition}} \\
\midrule
Positivist & Specify measurable targets for the shared life & Adopt practices with established effects & Administer the instruments; compute the proxies & Do not possess the measure as the truth & Reset targets to the new baseline \\
Evidentialist & Form belief about the good in proportion to evidence & Employ supported practices within their limits & Update belief on the lived evidence & Revise belief; resist over-confidence in the measure & Begin the turn on the revised belief \\
Normative & Specify what this shared life ought to value & Make room for each party's capability to unfold & Test return and unfolding against the specification & Reassert the question of justice & Re-specify the good for the next turn \\
Virtue-theoretic & Discern the excellences this shared life requires & Cultivate, do not legislate, the dispositions (\xref{sec:prac-constraint}) & Read the dispositions formed & Guard against the hardening of habit into rule & Carry the cultivated disposition into the turn \\
Phenomenological & Attend to the lived sense of the good for these two & Cultivate the quality of attention & Describe thickly; refuse the metric's last word & Return the question to the unsymbolised register & Reopen the apprehension for the next turn \\
\bottomrule
\end{tabularx}
\end{table}

The two tables together state, in compact form, the foundational structure the paper develops: that the cognition and the practice of a shared flourishing are conducted at five moments, that each moment draws upon five frameworks, and that each framework contributes, at each moment, both a way of cognising the moment's object and a way of guiding its action. The prose of the sections that follow develops the cells of this matrix where each is most consequential, the normative and the virtue-theoretic at the moment of cognition (\xref{sec:cognition}), the virtue-theoretic and the evidential at the moment of practice (\xref{sec:practice}), the positivist most fully at the moment of evaluation (\xref{sec:evaluation}), and the reassertion of the unsymbolised and the just at the moments of reflection and re-cognition (\xref{sec:reflection}, \xref{sec:recognition}). The matrix does not replace that development, since the prose alone can state the tensions among the frameworks and the qualifications each entry requires; it provides the map within which the development is to be read, and the assurance that the treatment of the moments, though it foregrounds the framework most consequential at each, rests on the full matrix of frameworks and moments that this section sets out.
